\documentclass[book.tex]{subfiles}

\begin{document}

\chapter{Feynman-Kac Formula}

\label{chap:perturbation_theory}
\noindent\rule{11cm}{0.4pt} \\
\textbf{Prerequisites:} \autoref{chap:schrodinger} \\
\textbf{Difficulty Level:} ** \\
\noindent\rule{11cm}{0.4pt}
\vspace{5mm}
%\textcolor{red}{TODO: Cite Wipf path integrals}
%\url{https://www.scribd.com/document/525224183/A-Wipf-Path-Integrals}
The Feynman-Kac formula provides a powerful bridge between solutions to certain classes of PDEs and solutions to stochastic differential equations. In practice, the formula can provide a powerful formula to resolve certain classes of intractable PDEs.

\section{Statement}

Suppose that we have a parabolic PDE given by

\begin{align}
\frac{\partial u}{\partial t}(x, t) + \mu(x, t)\frac{\partial u}{\partial x}(x, t) + \frac{1}{2}\sigma^2(x, t)\frac{\partial^2 u}{\partial x^2}(x, t) - V(x, t)u(x, t) + f(x, t) &= 0
\end{align}
where $\mu, \sigma, V, f$ are known functions. Here we also impose a boundary condition that
\begin{align}
u(x, T) &= \psi(x)
\end{align}
where $\psi$ is another known function. The Feynman-Kac formula states that the solution to this PDE is given by the conditional expectiation

\begin{align}
u(x, t) &= E\left [ \int_t^T e^{-\int_t^r V(X_\tau, \tau) d \tau } f(X_\tau, r) dr + e^{- \int_t^T V(X_\tau, \tau) d \tau \psi(X_T)}  \right ]
\end{align}
Here, we have that $X_\tau$ is an It\^{o} process given by the equation
\begin{align}
d X_t \mu(X, T) dt + \sigma(X, t) dW_t
\end{align}
where $W$ is Wiener noise.
\section{Derivation}

Let us start by considering one particular instantiation of this formula and performing a derivation. Consider $K$ to be the poropagator of Schrodinger's equation. Note that Schrodinger's equation is a parabolic PDE, so we can use the Feynman-Kac formula. We start by using the Trotter formula
\begin{align}
    e^{A+B} &= \lim_{n\to\infty} (e^{A/n}e^{B/n})^n
\end{align}
%\textcolor{red}{TODO: Wipf uses a proof with a telescoping expansion. Perhaps re-use or use an alternative framing}
Here $A$ and $B$ are matrices. The Trotter product formula can be extended to the case $A$ and $B$ are self-adjoint operators and $A+B$ is self-adjoint on the intersection of the domains of $A$ and $B$.

\begin{align}
    e^{-it(A+B)} &= s - \lim_{n\to\infty} \left (e^{-itA/n}e^{-itB/n}\right )^n
\end{align}
Consider the unitary time evolution kernel
\begin{align}
    K(t, q, q') &= \langle q | e^{-itH}|q'\rangle 
\end{align}
Suppose that the Hamiltonian has form $H_0 + V$. Then
we split up the evolution into $n$ steps  by using the Trotter formula.
\begin{align}
    K(t, q, q') &= \lim_{n\to\infty} \langle q | \left (e^{-itH_0}e^{-itV} \right )^n|q'\rangle 
\end{align}
We insert the decomposition of identity given below $n$ times between the factors
\begin{align}
    I &= \int dw_j |w_j\rangle \langle w_j |
\end{align}
And we obtain
\begin{align}
    K(t,q,q') &= \lim_{n\to\infty} \int dw_1 \dotsc dw_n \prod_{j=0}^{n-1} \langle w_{n+1} | e^{-itH_0/n}e^{-itV/n}|w_j\rangle
\end{align}
We note that the potential is diagonal in the coordinate representation since it directly multiplies by a constant value at each point in position space. Thus we can write
\begin{align}
    \langle w_{n+1} | e^{-itH_0/n}e^{-itV/n}|w_j\rangle &= \langle w_{n+1} | e^{-itH_0/n}|w_j\rangle e^{-itV(w_j)/n}
\end{align}
Note that the free propagator is given by
\begin{align}
    K_0(t, q, q') &= \langle q' | e^{-itH} q \rangle \\
    &= \left(\sqrt{\frac{m}{2\pi i \hbar t}}\right)^d e^{im(q-q')^2/2\hbar t} \\
    &= A_t^d e^{im(q-q')^2/2\hbar t}
\end{align}
where
\begin{align}
    A_t &= \sqrt{\frac{m}{2\pi i \hbar t}}
\end{align}
Plugging this all together we obtain
\begin{align}
    K(t, q, q') &= \lim_{n\to\infty} A_{\epsilon}^n \int dw_1 \dotsc dw_{n-1}e^{iS^{(n)}(w)/\hbar} \\
    S^{(n)}(w) &= \frac{m}{2}\sum_{j=0}^{n-1}\epsilon \left ( \frac{w_{j+1} - w_j}{\epsilon}\right) - \sum_{j=0}^{n-1} \epsilon V(w_j)
\end{align}
This last quantity is a Riemann sum approximation to the action of a particle moving along the specified path
\begin{align}
    S^{(n)}(w) &\to \int_0^t ds (\frac{m}{2}\dot{w}^2 - V(w)) \\
    &= S(w)
\end{align}
This formula is often written thematically as
\begin{align}
K(t, q, q') &= \int_{w(0)=q}^{w(t) =q'} \mathcal{D}w e^{iS[w]/\hbar}
\end{align}
where the integral is taken over all possible paths.
%\textcolor{red}{TODO: Why is the $A_\epsilon^n$ quantity which diverges as $n\to\infty$ safe? Wipf says its fine by why?}
\end{document}